\chapter{Digital Ocean Docker Deployment}
\ChapterLabel{DigitalOceanDocker}
\index{Digital Ocean}
\index{Docker}
\index{Docker!deployment}
\index{UML class diagram}

\section{Overview}
This chapter records how a Digital Ocean Droplet was created, how Docker was
installed on it, and how a small Node.js website (the \emph{Color Buttons App})
was containerized and deployed. It also documents the refactor of the website's
JavaScript from inline event handlers to a class-based design, with a UML class
diagram of the updated design.

The deployed website is available at:
\begin{itemize}
  \item \url{http://144.126.196.140} (Droplet public IPv4 address)
\end{itemize}

\subsection{Team Contributions}
\begin{tabularx}{\linewidth}{L{4.0cm}Y}
\toprule \textbf{Member} & \textbf{Contribution}\\ \midrule
Rifat Rahman Khan (RRK) & Installed and ran everything: created the Droplet, installed Docker, and deployed the website, capturing the terminal sessions used as evidence.\\
Deanna Gaber (DMG) & Created the documentation for this report: wrote this chapter, the class-based JavaScript refactor, and the UML class diagram.\\ \bottomrule
\end{tabularx}

\section{Droplet Creation}
The Droplet was created in the Digital Ocean control panel with the following
configuration (confirmed by the console screenshot in
Figure~\ref{fig:droplet}):

\begin{itemize}
  \item Name: \texttt{ubuntu-s-1vcpu-512mb-10gb-lon1}
  \item Image: Ubuntu 24.04 (LTS) x64
  \item Plan: 1 vCPU, 512 MB RAM, 10 GB disk (\$4.00/month)
  \item Region: London (LON1), project \texttt{first-project}
  \item Authentication: SSH key (\texttt{id\_ed25519})
  \item Public IPv4 address: \texttt{144.126.196.140}
\end{itemize}

\begin{figure}[H]
  \centering
  \includegraphics[width=0.9\textwidth]{figures/droplet-overview.png}
  \caption{Digital Ocean console showing the active Droplet and its public IPv4 address.}
  \label{fig:droplet}
\end{figure}

\section{Installing Docker on the Droplet}
After connecting to the Droplet over SSH as \texttt{root}, Docker was installed
from Ubuntu's package repositories and verified with the \texttt{hello-world}
image. The terminal session was recorded with \texttt{script}.

\begin{listing}[H]
\begin{minted}[fontsize=\small,breaklines,frame=lines]{bash}
# Connect to the Droplet
ssh -i ~/.ssh/id_ed25519 root@144.126.196.140

# Install Docker (Ubuntu 24.04 package: docker.io)
apt update
apt install docker.io -y
systemctl enable --now docker

# Verify the installation
docker --version
docker run hello-world
\end{minted}
\caption{Installing and verifying Docker on the Droplet.}
\end{listing}

\noindent The resulting version was \texttt{Docker version 29.1.3, build
29.1.3-0ubuntu3\textasciitilde 24.04.2}, and \texttt{hello-world} printed
``Hello from Docker!'', confirming the daemon works.

\section{Smoke Test with nginx}
Before deploying our own application, a stock nginx container was started on
port 80 to prove the Droplet was reachable from the internet. It was then
stopped and removed so the port could be reused.

\begin{listing}[H]
\begin{minted}[fontsize=\small,breaklines,frame=lines]{bash}
docker run -d --name website -p 80:80 nginx
docker ps
docker stop website
docker rm website
docker ps
\end{minted}
\caption{nginx smoke test: start, verify, stop, and remove.}
\end{listing}

\begin{figure}[H]
  \centering
  \includegraphics[width=0.8\textwidth]{figures/nginx-welcome.png}
  \caption{The nginx welcome page served from \url{http://144.126.196.140}.}
  \label{fig:nginx}
\end{figure}

\section{Deploying the Color Buttons App}
The application consists of a small Node.js server (\texttt{server.js}), a
\texttt{package.json}, a static \texttt{public/} directory containing
\texttt{index.html}, and a \texttt{Dockerfile}.

\subsection{Dockerfile}
\begin{listing}[H]
\inputminted[fontsize=\small,breaklines,frame=lines]{docker}{color-buttons-app/Dockerfile}
\caption{Dockerfile for the Color Buttons App.}
\end{listing}

\subsection{Copy, Build, and Run}
The project directory was copied from the local machine to the Droplet with
\texttt{scp}, then built into an image and run as a detached container with the
container's port 3000 published on the Droplet's port 80.

\begin{listing}[H]
\begin{minted}[fontsize=\small,breaklines,frame=lines]{bash}
# On the local machine: copy the app to the Droplet
scp -i ~/.ssh/id_ed25519 -r color-buttons-app root@144.126.196.140:/root/

# On the Droplet
ssh -i ~/.ssh/id_ed25519 root@144.126.196.140
cd /root/color-buttons-app
docker build -t color-buttons-app .
docker run -d --name color-buttons-app -p 80:3000 color-buttons-app

# Verify
docker ps
curl http://localhost
\end{minted}
\caption{Commands used to deploy the application.}
\end{listing}

The build succeeded (\texttt{Successfully tagged color-buttons-app:latest}),
\texttt{docker ps} showed the container \texttt{Up} with
\texttt{0.0.0.0:80->3000/tcp}, and \texttt{curl http://localhost} returned the
page's HTML.

\begin{figure}[H]
  \centering
  \includegraphics[width=0.8\textwidth]{figures/color-buttons.png}
  \caption{The deployed Color Buttons App at \url{http://144.126.196.140}.}
  \label{fig:buttons}
\end{figure}

\subsection{Environment Checks}
The following commands were also run to record the environment.

\begin{listing}[H]
\begin{minted}[fontsize=\small,breaklines,frame=lines]{bash}
cat /etc/os-release   # Ubuntu 24.04.5 LTS (Noble Numbat)
ufw status            # Status: inactive
\end{minted}
\caption{Environment checks.}
\end{listing}

\noindent The firewall (\texttt{ufw}) was \emph{inactive} for this exercise, so
port 80 was reachable without further rules. A production deployment should
enable \texttt{ufw} (allowing only SSH and HTTP/HTTPS) and serve the site over
HTTPS; the browser currently labels the site ``Not Secure'' because it uses
plain HTTP.

\section{Class-Based JavaScript Refactor}
The original page registered two anonymous arrow functions directly on the
\texttt{blueBtn} and \texttt{redBtn} elements, duplicating the same logic twice.
The updated \texttt{public/index.html} replaces this with two classes:

\begin{itemize}
  \item \texttt{ColorButton} wraps a single \texttt{<button>} element and the
        color it applies. It exposes \texttt{apply()} to set the background color and
        \texttt{bind()} to attach the click listener.
  \item \texttt{ColorButtonApp} owns a collection of \texttt{ColorButton}
        objects, builds them from a configuration list, and starts them with
        \texttt{start()}.
\end{itemize}

Adding a new color now requires only one new entry in the configuration list
rather than new copy-pasted handler code.

\begin{listing}[H]
\begin{minted}[fontsize=\small,breaklines,frame=lines,linenos]{javascript}
class ColorButton {
  #button;
  #color;
  #target;

  constructor(buttonId, color, target = document.body) {
    this.#button = document.getElementById(buttonId);
    this.#color = color;
    this.#target = target;
  }

  get color() {
    return this.#color;
  }

  apply() {
    this.#target.style.backgroundColor = this.#color;
  }

  bind() {
    this.#button.addEventListener("click", () => this.apply());
  }
}

class ColorButtonApp {
  #buttons = [];

  constructor(configs) {
    this.#buttons = configs.map(c => new ColorButton(c.id, c.color));
  }

  start() {
    this.#buttons.forEach(button => button.bind());
  }
}

new ColorButtonApp([
  { id: "blueBtn", color: "blue" },
  { id: "redBtn",  color: "red"  }
]).start();
\end{minted}
\caption{Class-based JavaScript in \texttt{public/index.html}.}
\end{listing}

\subsection{Redeploying the Updated Site}
After editing \texttt{index.html}, the container image must be rebuilt and the
old container replaced:

\begin{listing}[H]
\begin{minted}[fontsize=\small,breaklines,frame=lines]{bash}
# Copy the updated page to the Droplet
scp -i ~/.ssh/id_ed25519 color-buttons-app/public/index.html \
    root@144.126.196.140:/root/color-buttons-app/public/

# Rebuild the image and replace the running container
ssh -i ~/.ssh/id_ed25519 root@144.126.196.140 '
  cd /root/color-buttons-app &&
  docker build -t color-buttons-app . &&
  docker stop color-buttons-app &&
  docker rm color-buttons-app &&
  docker run -d --name color-buttons-app -p 80:3000 color-buttons-app'
\end{minted}
\caption{Redeploying after the class-based change.}
\end{listing}

\subsection{UML Class Diagram}
Figure~\ref{fig:uml} shows the updated design. \texttt{ColorButtonApp} is
composed of one or more \texttt{ColorButton} objects, and each
\texttt{ColorButton} depends on the browser's \texttt{HTMLElement} (the button
and the target element whose background changes). Private members are marked
with \texttt{-} and public members with \texttt{+}.

\begin{figure}[H]
\centering
\begin{tikzpicture}[
  font=\small,
  class/.style={rectangle split, rectangle split parts=3, draw, thick,
                text width=6.4cm, align=left, fill=gray!8},
  ext/.style={rectangle split, rectangle split parts=1, draw, thick,
              text width=3.2cm, align=center, fill=blue!6}
]
  \node[class] (app) at (0,0) {
    \textbf{\hfill ColorButtonApp \hfill\null}
    \nodepart{second}
    -\,buttons: ColorButton[]
    \nodepart{third}
    +\,constructor(configs: Config[])\\
    +\,start(): void
  };
  \node[class] (btn) at (0,-4.6) {
    \textbf{\hfill ColorButton \hfill\null}
    \nodepart{second}
    -\,button: HTMLElement\\
    -\,color: string\\
    -\,target: HTMLElement
    \nodepart{third}
    +\,constructor(buttonId, color, target)\\
    +\,color: string \{get\}\\
    +\,apply(): void\\
    +\,bind(): void
  };
  \node[ext] (el) at (7,-4.6) {\textbf{HTMLElement}\\{\scriptsize (browser DOM)}};

  % composition: App owns 1..* ColorButton
  \draw[thick] (app.south) -- node[right,pos=.45]{\scriptsize 1..*} (btn.north);
  \node[draw,fill=black,diamond,inner sep=2pt] at ($(app.south)+(0,-0.15)$) {};
  % dependency
  \draw[thick,dashed,-{Stealth}] (btn.east) -- node[above]{\scriptsize uses} (el.west);
\end{tikzpicture}
\caption{UML class diagram of the class-based Color Buttons design.}
\label{fig:uml}
\end{figure}

\section{Retrospective}
The Droplet, Docker installation, nginx smoke test, and application
deployment all worked on the first attempt. Two follow-ups are worth noting:
the firewall should be enabled and HTTPS added before this is used beyond a
class exercise, and Docker's legacy builder printed a deprecation warning
(installing \texttt{docker-buildx} would remove it).
